\chapter{Conclusiones}

El desarrollo de este trabajo ha permitido estudiar FlashAttention desde dos niveles de abstracción muy diferentes. Por una parte, se ha implementado una versión en Triton, delegando una parte importante de las decisiones de bajo nivel al compilador. Por otra, se ha desarrollado directamente en CUDA un kernel específico para arquitecturas Ampere, controlando de forma explícita la distribución del trabajo entre \textit{warps}, el movimiento de datos entre memoria global, memoria compartida y registros, las instrucciones \texttt{cp.async} y \texttt{ldmatrix}, el uso de Tensor Cores mediante \texttt{mma.sync} y la disposición de los datos en memoria compartida.

El desarrollo de la implementación CUDA ha requerido un esfuerzo considerablemente mayor. Durante este proceso se han identificado y corregido problemas relacionados con accesos globales no coalescentes, conflictos de banco, presión sobre el fichero de registros, disposición de los fragmentos utilizados por las instrucciones matriciales, sincronización entre \textit{warps} y utilización de múltiples buffers en memoria compartida. Los perfiles finales muestran que muchas de estas optimizaciones alcanzaron el objetivo previsto: los accesos globales son coalescentes, no se producen conflictos de banco en la implementación CUDA, no existe \textit{register spilling} y el consumo de registros y memoria compartida resulta incluso inferior al de algunas de las implementaciones de referencia.

Sin embargo, los resultados muestran también una de las principales lecciones obtenidas durante el trabajo: optimizar correctamente los detalles microarquitectónicos de un kernel no garantiza que la estrategia global de distribución del trabajo sea la adecuada.

La implementación CUDA procesa bloques relativamente pequeños de consultas y presenta una reutilización de \(K\) y \(V\) inferior a la observada directamente en el \textit{mapping} generado por Triton. Al mismo tiempo, para determinados tamaños del problema genera una cantidad limitada de trabajo independiente, impidiendo aprovechar completamente los recursos que teóricamente permitirían mantener varios CTAs residentes en cada SM. En cuDNN, el menor número de sectores solicitado es consistente con una mayor reutilización efectiva, pero el kernel interno no se reconstruye y por tanto no se atribuye esa diferencia a un tamaño de tile concreto. El perfil de SDPA muestra, por su parte, un número de bloques significativamente superior.

Durante el desarrollo preliminar se realizaron también pruebas informales sobre una RTX 3060 Laptop en las que el comportamiento relativo de la implementación CUDA parecía más favorable. Estas ejecuciones no se realizaron bajo el protocolo experimental final, no se documentaron con el mismo entorno, versiones, repeticiones y perfiles, y por tanto \textbf{no forman parte de la evaluación experimental ni se extrae de ellas ninguna conclusión}. Se conservan únicamente como una observación anecdótica que motivó la posterior evaluación controlada sobre la A100.

Las medidas realizadas para longitudes de secuencia crecientes permiten además matizar los resultados. A medida que aumenta \(N\), la implementación CUDA dispone de un número mayor de CTAs y la diferencia respecto a las alternativas se reduce. Para las secuencias grandes ensayadas, el kernel desarrollado alcanza aproximadamente \(100\) TFLOP/s sostenidos, frente a aproximadamente \(126\) TFLOP/s de Triton, \(167\) TFLOP/s de cuDNN y \(177\) TFLOP/s de SDPA. Estas cifras permiten valorar el resultado de forma cuantitativa y sin recurrir a calificativos subjetivos.

En este sentido, el resultado del proyecto no debe medirse únicamente por cuál de las implementaciones obtiene el menor tiempo de ejecución. El desarrollo manual ha permitido comprender aspectos que permanecen ocultos al utilizar bibliotecas de alto nivel: cómo se distribuyen los fragmentos de una instrucción \texttt{mma.sync} entre los hilos de un \textit{warp}, cómo se pueden utilizar transferencias asíncronas para construir un \textit{pipeline}, cómo afecta el patrón de acceso a los bancos de memoria compartida, cómo interactúan registros, memoria compartida y ocupación, y por qué una mayor ocupación teórica no implica necesariamente un mayor rendimiento.

Asimismo, la comparación con Triton ha resultado especialmente útil. Triton permite expresar el mismo algoritmo con una cantidad de código considerablemente menor y alcanza aproximadamente \(126\) TFLOP/s en el régimen de secuencias grandes medido, lo que pone de manifiesto el valor de los compiladores especializados para programación GPU. Sin embargo, el desarrollo CUDA proporciona una visibilidad mucho mayor sobre las decisiones que finalmente determinan el rendimiento y permite interpretar con mayor profundidad el código generado por herramientas de más alto nivel.

\section{Limitaciones}

Los resultados de este TFG deben interpretarse dentro de un alcance experimental deliberadamente acotado. La evaluación principal se realiza sobre una única NVIDIA A100-SXM4-40GB y, por tanto, no demuestra que las mismas relaciones de rendimiento se mantengan en otras GPUs Ampere, en generaciones posteriores o en aceleradores de otros fabricantes. Además, la instancia se obtiene a través de Vast.ai y las frecuencias no se fijan manualmente; las desviaciones típicas observadas son pequeñas, pero no se dispone de un entorno de frecuencia completamente controlada.

La configuración estudiada utiliza BF16, dimensión de cabeza fija \(d=128\), \(B=H=1\), atención no causal y longitudes de secuencia múltiplos de 16. No se evalúan otras dimensiones de cabeza, varios batches o cabezas simultáneos, atención causal, dropout, secuencias arbitrarias ni formatos numéricos distintos. Tampoco se estudian backward, entrenamiento completo, inferencia autoregresiva con caché KV ni integración dentro de un modelo completo; las conclusiones se refieren exclusivamente al forward aislado analizado.

El perfilado microarquitectónico detallado con Nsight Compute corresponde a \(N=8192\). Las medidas temporales cubren un rango mucho mayor de longitudes, pero no debe asumirse que porcentajes concretos de ocupación, actividad Tensor Core o comportamiento de caché sean constantes para todos esos tamaños. Del mismo modo, el análisis de cuDNN se limita a las métricas observables del kernel generado; no se reconstruye su implementación interna y, por ello, las inferencias sobre reutilización se formulan como interpretaciones consistentes con los contadores, no como hechos sobre su \textit{mapping}.

Finalmente, la validación funcional se basa en \texttt{torch.allclose} con una semilla fija y tolerancias de \(2\cdot10^{-2}\). Esto es suficiente como comprobación de regresión para el desarrollo realizado, pero no constituye un estudio exhaustivo de estabilidad numérica. Una evaluación dedicada a precisión debería utilizar múltiples semillas, más familias de entradas y métricas de error máximo, medio y relativo.

\section{Trabajo futuro}

Como líneas de trabajo futuras, las principales modificaciones identificadas consisten en aumentar el tamaño del bloque de consultas procesado por CTA, estudiar el almacenamiento de \(Q\) en memoria compartida para incrementar la reutilización de \(K\) y \(V\), reducir la frecuencia de las actualizaciones del \textit{softmax online} y de las sincronizaciones, y explorar estrategias de partición que permitan generar mayor paralelismo global. Estas modificaciones implicarían cambios estructurales importantes sobre el kernel actual, pero atacan directamente las limitaciones observadas durante el análisis de los perfiles.

En definitiva, el trabajo confirma que FlashAttention es un problema especialmente adecuado para estudiar la programación de GPUs modernas, ya que obliga a considerar simultáneamente cómputo matricial, jerarquía de memoria, paralelismo, sincronización y utilización de recursos. Aunque el resultado final de la implementación CUDA no haya alcanzado el rendimiento esperado sobre la A100, el proceso de llegar hasta ese resultado ha permitido identificar con precisión por qué ocurre y qué decisiones deberían modificarse para continuar mejorándolo. Esa comprensión constituye, probablemente, uno de los resultados más importantes obtenidos durante el desarrollo del proyecto.
